\input{preamble}

\begin{document}

\title{Smooth Manifolds}
\maketitle

\phantomsection
\label{section-phantom}
\tableofcontents

\section{Orientable manifolds}
\label{section-orientable-manifolds}

\begin{exercise}
\label{exercise-orientable-surfaces}
Show that a regular surface $S \subset
\mathbb{R}^3$
is orientable if and only if there is a
differentiable
map $N:S \to \mathbb{R}^3$ such that
$N(p)\perp T_pS$
and $|N(p)|=1$ for all $p \in S$.
\end{exercise}

\begin{proof}
($\implies$) Take the normalization of the
cross product of an orientable basis.
It is nowhere vanishing because the
determinant,
which is positive by orientability,
is the size of the cross product.

($\impliedby$) Take the determinant of the
matrix $(v,w,N)$
where $v,w$ is any basis of $S$.
This is nowhere zero differentiable 2-form on
$S$.
\end{proof}

\section{Immersions and submersions}
\label{section-immersions-and-submersions}

\noindent
For any smooth manifold submersion
$\pi:\tilde{M}\to M$ there is a bundle
isomorphism I love:
\begin{equation}
\label{equation-subme-bundl-decom}
\pi^*TM \oplus \Ker \pi\cong
T\tilde{M}
\end{equation}
 where $\Ker\pi$ is the {\it kernel bundle},
 whose
fiber is the kernel of $\pi_*$, and is also
called the {\it vertical bundle}.
Proving that $\Ker\pi$ is in fact a bundle
and confirming the isomorphism
\ref{equation-subme-bundl-decom} is a simple
exercise.

We call $\pi$ a {\it riemannian submersion}
only when $\pi^*TM$ is isometric to
$TM$, which makes sense only when $\tilde{M}$
and $M$ are Riemannian manifolds.

Following notation of
\cite{Milnor-Characteristic-Classes}, for any
subbundle
$\xi \subset\eta$ over a Riemannian manifold,
the {\it orthogonal complement
bundle} $\xi^\perp$ can be defined fiberwise
as the orthogonal complement of the
fiber and shown to be a bundle as in the case
of the kernel bundle. Then Milnor
defines the {\it normal bundle} when $\xi$ is
the tangent bundle of a
submanifold. Finally this construction is
noted to work just as well for any
isometric immersion $f:M\to \tilde{M}$, so
that
\begin{equation}
\label{equation-immer-bundl-decom}
f^*T\tilde{M}\cong TM\oplus T^\perp M
\end{equation}
Further, in problem 3-B we are asked to
define for $\xi \subset \eta$ the {\it
quotient bundle} and show it's a bundle. And
if $\eta$ has a Riemannian metric,
then $\xi^\perp\cong\eta/\xi$. So that's why
we sometimes say that we can define
the normal bundle using a Riemannian metric
but we don't need to.

An interesting thought is: given a subbundle,
is there a manifold whose tangent
bundle is that subbundle?

\section{Embeddings}
\label{section-embeddings}

\begin{exercise}
\label{exercise-p2-in-r4}
Seja \(F:\mathbb{R}^3 \to \mathbb{R}^4\)
dada por
\[F(x,y,z)=(x^2-y^2,xy,xz,yz),\qquad
(x,y,z)=p \in \mathbb{R}^3.\]
Seja \(S^2 \subset \mathbb{R}^3\) a esfera
unitária com centro na origem \(0 \in
\mathbb{R}^3\). Oberve que a restrição \(
\varphi:= F|_{S^2}\) é tal que
\(\varphi(p)=\varphi(-p)\), e considere a
aplicação \(\tilde{\varphi}:\mathbb{R}P^2
\to
\mathbb{R}^4\) dada por
\[\tilde{\varphi}([p])=\varphi(p),\qquad
[p]\text{=clase de equivalência de
\(p=\{p,-p\}\)}\]

Prove que
\begin{enumerate}
\item \(\tilde{\varphi}\) é uma imersão.
\item \(\tilde{\varphi}\) é biunívoca;
junto com (a) e a compacidade de
\(\mathbb{R}P^2\), isto implica
que \(\tilde{\varphi}\) é um mergulho.
\end{enumerate}
\end{exercise}

\begin{proof}
\begin{enumerate}
\item Considere a carta \(\{z=1\}\). A
representação coordenada de
\(\tilde{\varphi}\) vira
\[(x,y) \longmapsto (x^2-y^2,xy,x,y)\]
cuja derivada como mapa \(\mathbb{R}^2 \to
\mathbb{R}^4\) é
\[\begin{pmatrix} 2x& -2y\\y&x\\1&0\\0&1
\end{pmatrix}\]
que é injetiva. Agora pegue a carta
\(\{x=1\}\). Então a representão coordenada
de \(\tilde{\varphi}\) vira
\[(y,z) \longmapsto (1-y^2,y,z,yz)\]
e tem derivada
\[\begin{pmatrix} -2y&0\\1&0\\0&1\\z&y
\end{pmatrix}\]
que também é injetiva. Seguramente algo
análogo acontece na carta \(\{y=1\}\).

\item \(\tilde{\varphi}\) é injetiva. Pegue
dois pontos \(p_1:=[x_1:y_1:z_1]\) e
\(p_2:=[x_2:y_2:z_2]\) e suponha que
\(\tilde{\varphi}(p_1)=\tilde{\varphi}(p_2)\)
.
I.e.,
\[x_1^2-y_1^2=x_2^2-y_2^2,\qquad
x_1y_1=x_2y_2,\qquad x_1z_1=x_2z_2,\qquad
y_1z_1=y_2z_2\]
Suponha primeiro que \(z_1 \neq 0\). Segue
que
\[x_1=\frac{z_2}{z_1}x_2,\qquad
y_1=\frac{z_2}{z_1}y_2\]
logo
\[x_2^2-y_2^2=x_1^2-y_1^2=
\left(\frac{z_2}{z_1}\right)^2(x_2^2-y_2^2)
\implies
z_2=z_1\implies x_1=x_2,\qquad y_1=y_2\]

\noindent
Em fim, uma imersão injetiva com domínio
compacto é um mergulho porque é fechada:
pegue um fechado no domínio, vira compacto,
imagem é compacta, que é fechado. Pronto.
\end{enumerate}
\end{proof}

\begin{exercise}
\label{exercise-extension-vector-fields}
Dada uma subvariedade \(M \subseteq
\tilde{M}\) uma subvariedade mergulhada e
\(X \in \mathfrak{X}(M)\). Mostre que existe
um aberto \(U \subset \tilde{M}\) contendo
\(M\) e um campo
\(\tilde{X} \in \mathfrak{X}(U)\) tal que
\(\tilde{X}|_{M}=X\).
Caso \(M\) seja subconjunto fechado de
\(\tilde{M}\), prove que \(U\) pode ser
tomado igual a \(\tilde{M}\). Se \(M\) não
é subconjunto fechado de \(\tilde{M}\),
pode não existir extensão de \(X\)
definida em todo \(\tilde{M}\).
\end{exercise}

\begin{proof}
Take submanifold coordinates for $M$,
i.e. coordinates where $M\cap U
=(x_1,...,x_n,0,...,0)$.
The local expression of $X$ is
$\sum a^i(p)\partial_i$.
Then we just literally extend this to a
vector
field in all of $U$ by putting
$\tilde{X}=\sum
a^i(\pi_M(\tilde{p}))\partial_i$.
\end{proof}

\section{Principal bundles}
\label{section-principal-bundles}

Upshot:
a principal bundle is a fiber bundle such
that, if you choose a point in a fiber,
then the map $g \mapsto p\cdot g$
identifies that fiber with $G$.
It is surjective by transitivity and
injective by freeness.
The chosen point plays the role of the
identity.

\begin{definition}
\label{definition-principal-bundle}
Let $G$ be a Lie group.
A {\it principal $G$-bundle}
\[
\pi:P\to M
\]
is a smooth fiber bundle with a smooth
right action of $G$ on $P$ such that:
\begin{enumerate}
\item the action preserves fibers,
\item it is free and transitive on each
fiber, and
\item locally $P|_U \cong U\times G$
as $G$-spaces,
with $G$ acting on itself by right
multiplication.
\end{enumerate}
\end{definition}

Equivalently,
the transition functions of a local
trivialization take values in $G$.

In particular,
the frame bundle $\operatorname{Fr}(M)$
from
hodge-birational-atoms-2026.tex
is a principal
\(\operatorname{GL}(n,\mathbf R)\)-bundle.

\end{document}
